% ===========================================================================
%  01 — Introduction
%
%  Evidence: paper_context/proposal_summary.md
%            paper_context/system_delta_from_proposal.md
%            paper_context/claims_registry.md
%  Terminology: paper_context/terminology.md
% ===========================================================================
\chapter{Introduction}\label{ch:introduction}

% ===========================================================================
\section{Motivation and Problem Statement}\label{sec:intro-motivation}

High-fidelity flight simulators are increasingly capable of reproducing the
instrument density and operational tempo of real cockpit procedures. That
fidelity is valuable for training, but it also creates a practical tutoring
problem. A novice learner does not simply need to know the next checklist item.
They must recognise whether the current aircraft state makes that item valid,
verify that preconditions have been satisfied, notice when a display has reached
the intended state, and recover when an earlier action was skipped or only
partly completed. In this setting, procedural knowledge is not only textual; it
is distributed across switches, gauges, warning lights, display pages, and the
sequence history of the task.

Conventional aids such as manuals, videos, and printed checklists remain useful
for preparation, yet they are poorly matched to moments where the learner is
already inside the simulator and unsure about the current state. Looking away
from the cockpit breaks the training flow, and a generic text answer may be
fluent without being grounded in the aircraft state that caused the learner to
ask for help. The problem is therefore not whether a language model can describe
an F/A-18C procedure in general. The harder question is whether a tutoring
runtime can constrain generated assistance to what is actually evidenced in the
simulator at the moment of help.

This thesis studies that question through the F/A-18C cold-start procedure in
DCS World. The current procedure pack contains 33 ordered steps (S01--S33)
grouped into six operational phases, from initial electrical power through
engine start, avionics configuration, flight-control checks, and final pre-taxi
configuration (Section~\ref{sec:bg-coldstart}). The benchmark and VLM
evaluation reported in this thesis focus on the original 25-step core procedure;
the later extension to S33 adds further pre-taxi checks without changing the
central evidence problem. Some steps are directly observable through DCS-BIOS
telemetry, some depend on cockpit display content that must be interpreted from
images, and a small number remain manual or outside the instrument layout used
by the vision pipeline. This mixture makes the cold-start a compact but
representative case for state-dependent procedural tutoring.

The central premise of the thesis is that procedural assistance becomes more
reliable and more auditable when every actionable tutoring output is tied to an
explicit evidence source. In the system developed here, such evidence may be a
telemetry variable, a procedure-gate evaluation, a retrieved manual or checklist
snippet, or a VLM-derived visual fact. The language model is therefore not asked
to infer procedure state from its parametric knowledge alone. It operates inside
a runtime that supplies state evidence, checks structured outputs, validates
overlay targets, and records the resulting help cycle for replay and analysis.

% ===========================================================================
\section{Research Questions}\label{sec:intro-rq}

The thesis is guided by one main research question and three subordinate
questions:

\paragraph{Main RQ.}
How can an evidence-grounded multimodal tutoring system support
state-dependent procedural training in an immersive flight simulator?

\paragraph{RQ1 --- System architecture and evidence integration.}
How can simulator telemetry, retrieved procedural knowledge, and visual fact
observations be integrated so that tutoring outputs are state-aware,
auditable, and safe for procedural guidance?

\paragraph{RQ2 --- VLM adaptation and ontology design.}
To what extent can a small-data LoRA-adapted VLM reliably extract structured
cockpit visual facts, and how does visual fact ontology design affect critical
false positives?

\paragraph{RQ3 --- Formative VR pilot evidence.}
In a VR-based F/A-18C cold-start training task, what formative evidence does a
small pilot study provide about task completion, procedural errors, and tutor
interaction in a with-tutor versus without-tutor comparison?

RQ1 and RQ2 are addressed through the implemented runtime, replay and harness
infrastructure, and technical benchmark evaluation. RQ3 is addressed through a
completed formative VR pilot study ($n=9$). The pilot is treated as descriptive
evidence about feasibility, instrumentation, completion, and error patterns; it
is not used to claim controlled learning gains, workload reduction, retention,
or transfer.

% ===========================================================================
\section{Thesis Vision and Scope}\label{sec:intro-vision}

The original research proposal framed the project as a VR-first grounded tutor
for the Meta Quest~3 and DCS-VR~\citep{ThesisProposal1128}. In that vision, the
learner would interact with the cockpit through gaze, hand tracking, and voice;
the tutor would retrieve relevant manual and checklist passages; an LLM would
generate concise step guidance; and the result would be shown as in-headset
cards or spatial overlays. The proposal also described a three-condition human
evaluation comparing a video-and-notes baseline, an ungrounded text-only LLM,
and a grounded tutor, with outcome measures such as completion rate, procedural
errors, task time, NASA-TLX workload, knowledge quizzes, retention, and
transfer.

The completed thesis keeps the same long-term goal but sharpens the research
focus. Instead of treating VR presentation as the main contribution, it asks
what technical evidence chain is required before generated help can be trusted
in a live simulator procedure. This led to three implementation priorities that
were not fully visible in the proposal: a telemetry pipeline that distinguishes
known state from missing state, a visual-fact pipeline that adapts VLMs to
cockpit display evidence, and a structured help cycle that validates model
outputs against procedure gates, evidence references, and overlay constraints.
The VR path was implemented and exercised in a completed pilot study, but the
human evidence remains formative.

The thesis scope is therefore deliberately split into technical validation and
formative human evidence:

\begin{enumerate}
\item The design and implementation of an evidence-grounded tutoring runtime
  that combines DCS-BIOS telemetry, procedure-pack gates, retrieved procedural
  knowledge, VLM-derived visual facts, structured response validation, overlay
  display, logging, and replay (Chapters~\ref{ch:system-design},
  \ref{ch:methodology}, and~\ref{ch:implementation});
\item The technical validation of the visual fact adaptation pipeline through
  independent holdout benchmarks and cross-backbone comparisons
  (Chapter~\ref{ch:evaluation}, Part~A);
\item A completed formative VR pilot study ($n=9$) that provides descriptive
  evidence about task completion, error patterns, tutor interaction, and
  instrumentation readiness under with-tutor and without-tutor conditions
  (Chapter~\ref{ch:evaluation}, Part~B).
\end{enumerate}

% ===========================================================================
\section{Implemented Evidence Chain}\label{sec:intro-scope}

The implemented runtime can be understood as an evidence chain rather than as a
collection of isolated modules. At the bottom of the chain, DCS-BIOS telemetry
provides the most authoritative signal for cockpit controls, lights, and
numerical engine or system values. Raw UDP frames are normalised, sanitised, and
aggregated into stable runtime variables so that procedure gates can distinguish
between a false value and a missing source. This distinction matters because a
tutor should not advance a learner simply because an expected signal was absent.

Where telemetry is insufficient, the runtime uses a visual fact layer. Cockpit
display pages such as the FCS page, BIT root page, FCS-MC BIT states, and INS
alignment text are captured as composite-panel images and converted into
structured facts with states \texttt{seen}, \texttt{not\_seen}, or
\texttt{uncertain}. The VLM is restricted to this perception task. It does not
decide that the procedure is complete; downstream gates combine visual facts
with telemetry and sequence state.

The tutoring response is produced only after the runtime has assembled an
evidence packet containing the current procedure candidates, relevant telemetry
and visual observations, and retrieved procedural sources. The LLM/VLM output is
validated against runtime-generated JSON Schema constraints and mapped through
allowlisted overlay targets. If the model path is unavailable or invalid, the
system falls back to deterministic, text-only guidance. This design makes the
help cycle auditable: a later replay can inspect which evidence was available,
which step was inferred, what the model returned, and how the response was
validated before display.

Procedure definitions, gates, telemetry mappings, visual fact bindings, and
error taxonomies are kept in declarative procedure-pack configuration rather
than embedded in prompts. This pack-based design keeps the architecture aligned
with Clean / Hexagonal principles: domain rules remain separate from simulator
adapters, vision adapters, model adapters, and export tooling.

% ===========================================================================
\section{Contributions}\label{sec:intro-contributions}

This thesis makes three contributions.

\begin{enumerate}

\item \textbf{An evidence-grounded procedural tutoring architecture.}
The thesis presents a ports-and-adapters tutoring runtime in which each
actionable help output must be traceable to telemetry, a procedure-gate
evaluation, a retrieved source, or a visual fact. The contribution is not merely
that an LLM is connected to a simulator. It is that generated tutoring is placed
inside a constrained help cycle with declarative procedure packs,
schema-validated responses, allowlisted overlay targets, deterministic
fallbacks, and replayable logs. This contribution addresses RQ1 and is described
in Chapters~\ref{ch:system-design} and~\ref{ch:implementation}, with validation
evidence in Chapter~\ref{ch:evaluation}, Section~\ref{sec:eval-rq1}.
% Evidence: C-ARCH-01, C-ARCH-02, C-ARCH-03, C-PROC-01, C-PROC-02,
%          C-TELEM-01, C-RAG-01, C-RUNTIME-01

\item \textbf{A visual fact ontology and small-data VLM adaptation pipeline.}
The thesis develops a 13-fact visual ontology for F/A-18C cockpit display state,
evolved from an 8-fact baseline that exposed the weakness of asking a VLM to
recognise compound procedural states directly. The revised ontology decomposes
such targets into locally observable visual atoms and leaves procedural
interpretation to downstream rules. The associated data pipeline spans DCS
screenshot capture, AI pre-labelling, human review, bilingual SFT row
generation, LoRA fine-tuning, and independent holdout benchmarking. On the Qwen
backbone, a LoRA adapter trained on 928 bilingual rows reaches approximately
0.99 fact accuracy on two independent holdouts and reduces critical false
positives by 64--89\%; the same data recipe also improves the Gemma backbone.
This contribution addresses RQ2 and is described in
Chapters~\ref{ch:system-design} and~\ref{ch:methodology}, with evaluation in
Chapter~\ref{ch:evaluation}, Sections~\ref{sec:eval-qwen-13fact}
and~\ref{sec:eval-backbones}.
% Evidence: C-VLM-01 through C-VLM-07, C-DATA-01, C-TRAIN-01

\item \textbf{A VR procedural training evaluation framework and formative
pilot study.}
The thesis delivers the instrumentation required to study evidence-grounded
procedural tutoring with real participants: experiment export tooling, baseline
and tutor-condition configurations, session launch support, replayable logs,
semi-automated pre-scoring, error coding guides, and an experiment operation
manual. A formative VR pilot study ($n=9$) was conducted with four with-tutor
and five without-tutor participants. All four tutor-condition participants
completed the full 33-step procedure; none of the five without-tutor
participants did. The observed error patterns also differed: tutor-condition
errors were mainly assistance-coupled, while without-tutor errors were mainly
omissions. These results support feasibility and instrumentation claims, but not
controlled claims about learning effectiveness. This contribution addresses RQ3
and is documented in Chapter~\ref{ch:evaluation}, Section~\ref{sec:eval-pilot},
and Appendix~\ref{ch:appendix}.
% Evidence: artifacts/experiments/fa18c_quest3_simitutor_2026/
% Status: pilot completed with real participant data (n=9). Claims are
% descriptive/formative only --- no inferential statistics, no learning
% outcome validation, no retention/transfer evidence.

\end{enumerate}

Taken together, these contributions define the boundary of the thesis. The work
demonstrates that evidence-grounded tutoring can be engineered, benchmarked, and
exercised in a formative VR study for a demanding simulator procedure. It does
not claim that educational effectiveness has been established beyond that
formative evidence.

% ===========================================================================
\section{Thesis Structure}\label{sec:intro-structure}

\textbf{Chapter~\ref{ch:background}} explains why the F/A-18C cold-start task
requires multiple evidence sources, introduces the visual fact ontology, and
positions the thesis against related work in intelligent tutoring, simulator
instrumentation, state-aware assistance, VLM adaptation, and grounded AR/VR
systems.

\textbf{Chapter~\ref{ch:system-design}} presents the evidence-grounded
architecture, including the ports-and-adapters structure, runtime data flow,
procedure and gating model, telemetry processing, replay infrastructure, and VLM
component.

\textbf{Chapter~\ref{ch:methodology}} describes the methods used to create and
validate the evidence chain: telemetry grounding, knowledge retrieval and source
policy, visual fact data collection, ontology evolution, LoRA fine-tuning, and
holdout benchmark design.

\textbf{Chapter~\ref{ch:implementation}} documents the implementation choices
that make the architecture operational, including core domain logic, adapter
boundaries, structured help generation, overlay mapping, VLM integration, and
experiment tooling.

\textbf{Chapter~\ref{ch:evaluation}} evaluates the thesis claims in two parts:
technical evidence for RQ1 and RQ2, followed by the formative VR pilot evidence
for RQ3.

\textbf{Chapter~\ref{ch:discussion}} interprets the results, explains the
evolution from the original proposal, draws design lessons from the ontology
revision, and states the limitations and future work.

\textbf{Chapter~\ref{ch:conclusion}} summarises the completed work against the
research questions and restates the claim boundaries of the thesis.

\textbf{Appendix~\ref{ch:appendix}} contains supplementary figures, benchmark
tables, dataset statistics, hyperparameter configurations, CLI references, and
the 8-fact baseline benchmark.
